% GENERATED FILE. Edit canonical lessons, then run npm run generate:books.
% canonical-source-hash: fnv1a64:7406801e5cf4b8b4
% canonical-lessons: HI-C13-sir, HI-C13-haath

\chapter{Head and Hand}
\label{ch:hi-head-and-hand}
\hlchaptermodality{13}

\begin{chapteropening}
\textbf{By the end of this chapter:} \emph{I can name the head and the hand in Hindi and trace both back to the Sanskrit words they wore down from.}

Say \hi{हाथ}, then \hi{हस्त}, then the shape of the change between them: the -st- cluster simplifies and the vowel lengthens.
\end{chapteropening}

\section[\hi{सिर}]{\hi{सिर} (sir) — the head}

\label{lesson:HI-C13-sir}



\begin{quote}
A short, everyday word — and, once you look, a worn-down piece of a much older Sanskrit word you may recognize from yoga or anatomy.
\end{quote}

\subsection*{You'll want to know — The word}
\textbf{\hi{सिर}} (\emph{sir}) = \textquotedblleft{}\textbf{head}\textquotedblright{} — from Sanskrit \textbf{\hi{शिरस्}} (\emph{śiras}), \textquotedblleft{}head.\textquotedblright{} Centuries of everyday use wore the word down from the fuller, formal Sanskrit shape to the short, simple \emph{sir} you use every day.

Sanskrit's \emph{śiras} root is still alive in more formal/technical registers too — related \emph{śīrṣa}-based vocabulary shows up in yoga and Ayurvedic anatomical terms (as in poses and positions named for the head), even while everyday Hindi kept the shorter, worn-down \emph{sir}.

\subsection*{Your turn}
Say these aloud:

\begin{itemize}
\raggedright
  \item \textquotedblleft{}sir\textquotedblright{} — head
  \item the fuller Sanskrit source — \textquotedblleft{}śiras\textquotedblright{}
  \item worn down over time — \textquotedblleft{}śiras $\to$ sir\textquotedblright{}
\end{itemize}

\begin{tcolorbox}[breakable,colback=teal!4,colframe=teal!35!black,title={Before you move on}]
What does \textbf{\hi{सिर}} mean, and what Sanskrit word is it worn down from? (\textbf{\textquotedblleft{}Head\textquotedblright{}} — from \textbf{\hi{शिरस्}} \emph{śiras}.) Where else does this root show up, in a more formal register? (\textbf{Yoga/anatomical terms} built on the related \emph{śīrṣa} root.)
\end{tcolorbox}
\section[\hi{हाथ}]{\hi{हाथ} (hāth) — the hand}

\label{lesson:HI-C13-haath}



\begin{quote}
This word looks nothing like its Sanskrit ancestor at first glance — but the change between them follows a real, traceable pattern, not a random drift.
\end{quote}

\subsection*{You'll want to know — The word}
\textbf{\hi{हाथ}} (\emph{hāth}) = \textquotedblleft{}\textbf{hand}\textquotedblright{} — from Sanskrit \textbf{\hi{हस्त}} (\emph{hasta}), \textquotedblleft{}hand.\textquotedblright{}

\subsection*{You'll want to know — The honest sound-change story}
Getting from \emph{hasta} to \emph{hāth} isn't one clean step — it passed through \textbf{Prakrit} (the everyday spoken languages that stood between Sanskrit and modern Hindi) along the way: Sanskrit's \textbf{‑st‑} consonant cluster regularly simplified in Prakrit into a \textbf{geminate} (doubled), \textbf{aspirated} consonant (giving an intermediate form like \emph{hattha}). Later, as that doubled consonant simplified further down to modern Hindi's single \textbf{‑th‑}, the \textbf{vowel just before it lengthened} — \emph{a} becoming \textbf{ā} — as if to make up for the sound that was lost. That's the general shape of the change: \textbf{‑ st‑ $\to$ (Prakrit geminate aspirate) $\to$ single aspirate + a longer vowel} — a recognized pattern in how Sanskrit consonant clusters evolved into modern Hindi, not a one-off coincidence for this word alone.

\subsection*{Your turn}
Say these aloud:

\begin{itemize}
\raggedright
  \item \textquotedblleft{}hāth\textquotedblright{} — hand
  \item the Sanskrit source — \textquotedblleft{}hasta\textquotedblright{}
  \item the shape of the change — \textquotedblleft{}‑st‑ cluster $\to$ simplified consonant + longer vowel\textquotedblright{}
\end{itemize}

\begin{tcolorbox}[breakable,colback=teal!4,colframe=teal!35!black,title={Before you move on}]
What does \textbf{\hi{हाथ}} mean, and what Sanskrit word is it from? (\textbf{\textquotedblleft{}Hand\textquotedblright{}} — from \textbf{\hi{हस्त}} \emph{hasta}.) What's the general shape of the sound change between them? (\textbf{Sanskrit's ‑st‑ cluster simplified} through Prakrit into a \textbf{geminate} aspirated consonant (\emph{hattha}), which later simplified to a single consonant — and the vowel before it \textbf{lengthened} to compensate.)
\end{tcolorbox}
